Wir nehmen an, dass der Sprinter gleichmässig beschleunigt. Seine mittlere Geschwindigkeit ist dann $\bar v = v/2$, und er braucht für die Strecke $s$ die Zeit
\begin{align}
 t &= \frac{s}{\bar v} = \tSpF = \frac{2\cdot\sSp}{\vSp} = \tSp
\end{align}
Seine Beschleunigung ist
\begin{align}
 a &= \frac{v}{t} = \aSpF = \aSp
\end{align}
Die beschleunigende Kraft ist die Haftreibungskraft der Bahn auf die Schuhe. Sie kann höchstens $\mu_\mathrm{H}\,F_\mathrm{N} = \mu_\mathrm{H}\,m\,g$ betragen:
\begin{align}
 m\,a &\leq \mu_\mathrm{H}\,m\,g
\end{align}
Damit muss mindestens gelten
\begin{align}
 \mu_\mathrm{H} &= \muHF \\
   &= \frac{\left(\vSp\right)^2}{2\cdot\gE\cdot\sSp} \\
   &= \muH \approx \result{\muHP{0}{2}}
\end{align}
